% ============================================================
%  Discussion (final section of the State of the Art)
% ============================================================

\section{Discussion}
\label{sec:discussion}

The previous sections have reviewed the architectures available for building speech-based conversational systems, the turn-taking mechanisms that coordinate who speaks when, the VAD and VAP layers that underpin them, the small but growing literature on multi-party conversational AI, the spectrum of roles that an AI agent can play inside a group, and the work on group decision-making with its survival-task paradigm. Taken together, this body of work points to a few recurring gaps that are visible across more than one of these threads.

The first is a modality gap. Most of the work on AI agents in group conversation is text-based. The surveys by Gu et al.~\cite{gu2022mpc}, Zheng et al.~\cite{zheng2022polyadic}, and Houde et al.~\cite{houde2025koala} consider almost exclusively chat-style interactions, and the most recent study of AI moderation in a decision-making task, by Dubini~\cite{dubini2024multiuser}, also runs in chat. The systems that operate in speech are restricted to dyadic settings with two human participants and a gaze channel, such as ARI~\cite{addlesee2024ari} and Furhat~\cite{axelsson2025furhat}. The reason is technical: Addlesee et al.~\cite{addlesee2020asr} report that commercial diarisation services have error rates between 49\% and 67\% in multi-party settings, and that overlapping speech adds another 20\% to the word error rate. As a result, audio-only group conversation with more than two human participants is still largely unexplored.

The second gap concerns turn-taking in groups. In dyadic settings, predictive turn-taking based on Voice Activity Projection~\cite{ekstedt2022vap} has produced models that anticipate speaker transitions from acoustic features alone. Extending this to multi-party is an open problem: Malik et al.~\cite{malik2024vap} note that the output space of VAP grows exponentially with the number of speakers, and no current model handles more than two channels reliably. Audio-only addressee detection collapses to about 27\% accuracy with concurrent speakers~\cite{axelsson2025furhat}, and the end-to-end speech-to-speech models such as Moshi~\cite{defossez2024moshi}, LLaMA-Omni~\cite{fang2024llamaomni}, and Freeze-Omni~\cite{wang2024freezeomni} are all designed for one human at a time. A multi-party speech system built today therefore has to commit to a cascaded pipeline with an explicit floor-management protocol, and the design effort shifts up to the layers above the pipeline.

The third gap is on the role of the AI in the group. Koala~\cite{houde2025koala} reports that a proactive AI in a text group was described by participants as ``a pedantic student who wouldn't create space for others''. Gao et al.~\cite{gao2025facilitation} compare an AI facilitator to a human facilitator in peer-support groups and find that the human still wins on every emotional measure. Dubini~\cite{dubini2024multiuser} reports that an LLM moderator in a text-based Murder Mystery task has only a subtle effect on participation dynamics and limited effect on decision accuracy. The group decision-making literature is more encouraging: Hall and Watson~\cite{hall1970normative} show that explicit consensus-seeking instructions improve group performance on survival tasks, and Hémon et al.~\cite{hemon2024constructive} show that the same effect carries over to online groups working through videoconferencing. The combined picture is that structured guidance helps groups, but the right amount and form of structure for an AI moderator in audio-only multi-party settings is not yet established.

A fourth gap is on evaluation. The text-based work on multi-party AI typically measures user preference and subjective satisfaction~\cite{houde2025koala,zheng2022polyadic}. The robotic speech work measures addressee accuracy~\cite{addlesee2024ari,axelsson2025furhat}. Empathic Co-facilitator~\cite{adikari2022empathic} and Gao et al.~\cite{gao2025facilitation} measure emotional connection. None of these works measures group performance in the sense that the group decision-making literature does, with deviation from an expert ranking and synergy with respect to the best individual member. The survival-task paradigm provides such measures, but it has not so far been combined with AI-moderated speech-based interaction in a controlled empirical study.

In summary, AI moderation of group conversation is an active line of research, but it has so far stayed in text, or in dyadic speech with gaze, or in single-task chatbot prototypes. Speech-based multi-party AI moderation, with audio-only operation, more than two human participants, and a quantitative evaluation grounded in the group decision-making literature, is an open space. The chapters that follow describe one attempt to occupy it.
